
\documentclass{standalone}
\usepackage{pgfplots}
\pgfplotsset{compat=1.18}

\pgfmathsetmacro{\muVal}{0.2}       % Massenverhältnis mu (illustrativ)
\pgfmathsetmacro{\mmu}{1-\muVal}    % 1 - mu
\pgfmathsetmacro{\qStern}{-\muVal}  % Position Stern  auf der q1-Achse
\pgfmathsetmacro{\qPlanet}{\mmu}    % Position Planet auf der q1-Achse
\pgfmathsetmacro{\pOne}{0}          % p1 (Schnitt)
\pgfmathsetmacro{\pTwo}{0}          % p2 (Schnitt)

\begin{document}
\begin{tikzpicture}[
    line cap=round,
    line join=round,
    >=triangle 45,
    x=1cm, y=1cm,
]
\begin{axis}[
    declare function={
        % \x = q1, \y = q2, Schnitt q3 = 0; +1e-5 regularisiert die Singularitäten
        rS(\x,\y) = sqrt((\x + \muVal)^2  + \y^2 + 0.00001);
        rP(\x,\y) = sqrt((\x - \mmu)^2    + \y^2 + 0.00001);
        V(\x,\y)  = -\mmu/rS(\x,\y) - \muVal/rP(\x,\y) - 0.5*(\x^2 + \y^2);
        H(\x,\y)  = 0.5*((\pOne - \y)^2 + (\pTwo + \x)^2) + V(\x,\y);
    },
    width=8cm,
    height=6cm,
    view={135}{30},
    axis lines=center,
    ticks=none,
    xlabel=$q_1$,
    ylabel=$q_2$,
    zlabel={$H(q_1,q_2)$},
    xlabel style={above left},
    ylabel style={above right},
    zlabel style={above},
    zmin=-4, zmax=2,
]
% Fläche
    \addplot3[
        surf,
        domain=-2:2,
        domain y=-2:2,
        samples=30,
        samples y=30,
        restrict z to domain=-4:2,
        opacity=0.6,
    ] {H(x,y)};

\end{axis}
\end{tikzpicture}
\end{document}
